\chapter{Pengamat Keadaan}
\label{bab:obs}
\enchap{State Observers}

\begin{tujuan}
	Setelah mempelajari bab ini, Anda mampu menjelaskan perlunya estimasi keadaan, merancang pengamat prediktor dan pengamat \textit{current} orde penuh, merancang pengamat orde tereduksi, menerapkan prinsip pemisahan, dan menurunkan fungsi transfer kompensator ekuivalen.
\end{tujuan}

\section{Mengapa estimasi keadaan}
Umpan balik keadaan memerlukan seluruh vektor $\mathbf{x}$, padahal sensor biasanya hanya mengukur sebagian. Menurunkan pengukuran secara numerik ($\Delta\theta/T$) memperkuat derau dan kuantisasi, sedangkan sensor tambahan menambah biaya. Alternatifnya adalah mengestimasi $\hat{\mathbf{x}}$ dari model, pengukuran $y$, dan masukan $u$.

Estimator kalang terbuka, yaitu salinan model $\hat{\mathbf{x}}[k+1] = \Phi\hat{\mathbf{x}}[k] + \Gamma u[k]$, memiliki galat $\tilde{\mathbf{x}} = \mathbf{x} - \hat{\mathbf{x}}$ yang berevolusi menurut $\tilde{\mathbf{x}}[k+1] = \Phi\tilde{\mathbf{x}}[k]$. Galat ini hanya meluruh secepat plant dan tidak meluruh sama sekali untuk plant dengan integrator.

\section{Pengamat prediktor}
\term{Pengamat}{observer} mengoreksi model dengan galat keluaran:
\begin{equation}
	\hat{\mathbf{x}}[k+1] = \Phi\hat{\mathbf{x}}[k] + \Gamma u[k] + L\big(y[k] - C\hat{\mathbf{x}}[k]\big).
\end{equation}
Dinamika galatnya
\begin{equation}
	\tilde{\mathbf{x}}[k+1] = (\Phi - LC)\tilde{\mathbf{x}}[k],
\end{equation}
sehingga $L$ dipilih agar nilai eigen $\Phi - LC$ berada di lokasi yang diinginkan. Ini dapat dilakukan sembarang jika dan hanya jika $(\Phi, C)$ teramati. Karena nilai eigen $\Phi - LC$ sama dengan nilai eigen $\Phi^\top - C^\top L^\top$, perancangan $L$ adalah masalah penempatan kutub yang \textbf{dual} pada pasangan $(\Phi^\top, C^\top)$:
\begin{equation}
	L = \alpha_e(\Phi)\,W_o^{-1}\begin{bmatrix}0 & \cdots & 0 & 1\end{bmatrix}^\top.
\end{equation}

\begin{contoh}[Estimasi kecepatan dari posisi]
	Motor posisi dengan $C = [1\ \ 0]$. Kutub kontroler Bab~\ref{bab:pp} berjari-jari 0,756; kutub pengamat dipilih lebih cepat di $0{,}40$ dan $0{,}45$, sehingga $\alpha_e(z) = z^2 - 0{,}85z + 0{,}18$ dan $L = [0{,}9687\ \ 1{,}7035]^\top$. Persamaan yang dieksekusi setiap cacah, dengan $e_y = y[k] - \hat\theta[k]$:
	\begin{align*}
		\hat\theta[k+1] &= \hat\theta + 0{,}0906\,\hat\omega + 0{,}0094\,u + 0{,}969\,e_y,\\
		\hat\omega[k+1] &= 0{,}8187\,\hat\omega + 0{,}1813\,u + 1{,}704\,e_y.
	\end{align*}
	Simulasi dengan masukan yang berubah tiga kali, derau pengukuran posisi $\sigma = 0{,}005$ rad, dan estimasi awal yang sengaja salah (Gambar~\ref{fig:b12-obs}) memberi RMS galat kecepatan 0,010 rad/s, dibandingkan 0,093 rad/s untuk beda hingga $\Delta y/T$. Galat awal hilang dalam beberapa cacah.
\end{contoh}

\begin{figure}[H]
	\centering
	\begin{tikzpicture}
		\begin{axis}[skd,width=12cm,height=5.4cm,xmin=0,xmax=6,ymin=-1.3,ymax=1.3,xlabel={$t$ (s)},ylabel={$\omega$ (rad/s)},legend style={at={(0.5,1.02)},anchor=south},legend columns=3]
			\addplot[skdgray!70,const plot,thin] table[x=t,y=wfd] {\dat{p13_obs.dat}};
			\addplot[black,const plot] table[x=t,y=w] {\dat{p13_obs.dat}};
			\addplot[skdorange,const plot,thick] table[x=t,y=what] {\dat{p13_obs.dat}};
			\legend{beda hingga $\Delta y/T$, sebenarnya, estimasi pengamat}
		\end{axis}
	\end{tikzpicture}
	\caption{Estimasi kecepatan motor dari pengukuran posisi}
	\label{fig:b12-obs}
\end{figure}

\section{Pengamat \textit{current}}
Pengamat prediktor menghitung $\hat{\mathbf{x}}[k]$ dari pengukuran hingga $k-1$, sehingga ada tunda satu cacah. Pengamat \textit{current} memakai pengukuran terbaru:
\begin{align}
	\bar{\mathbf{x}}[k] &= \Phi\hat{\mathbf{x}}[k-1] + \Gamma u[k-1] && \text{(prediksi)},\\
	\hat{\mathbf{x}}[k] &= \bar{\mathbf{x}}[k] + L_c\big(y[k] - C\bar{\mathbf{x}}[k]\big) && \text{(koreksi)},
\end{align}
dengan dinamika galat $\tilde{\mathbf{x}}[k+1] = (\Phi - L_cC\Phi)\tilde{\mathbf{x}}[k]$. $L_c$ dihitung dengan Ackermann dengan $C\Phi$ menggantikan $C$; untuk kutub 0,40 dan 0,45 diperoleh $L_c = [0{,}7802\ \ 2{,}0807]^\top$. Pengamat \textit{current} memberi estimasi lebih segar dan cocok bila waktu komputasi jauh lebih kecil dari $T$; pengamat prediktor menyediakan seluruh periode cacah untuk komputasi. Struktur prediksi--koreksi pengamat \textit{current} sama dengan filter Kalman (Bab~\ref{bab:lqr}).

\section{Pengamat orde tereduksi}
Bila sebagian keadaan terukur langsung, cukup estimasi sisanya. Partisi $\mathbf{x} = [x_a\ \ \mathbf{x}_b]^\top$ dengan $x_a = y$ terukur, dan matriks $\Phi$, $\Gamma$ dipartisi sesuai. Pengamat orde tereduksi adalah
\begin{equation}
	\hat{\mathbf{x}}_b[k+1] = \Phi_{bb}\hat{\mathbf{x}}_b + \Phi_{ba}y + \Gamma_bu + L_r\big(y[k+1] - \Phi_{aa}y - \Gamma_au - \Phi_{ab}\hat{\mathbf{x}}_b\big),
\end{equation}
dengan dinamika galat $\tilde{\mathbf{x}}_b[k+1] = (\Phi_{bb} - L_r\Phi_{ab})\tilde{\mathbf{x}}_b[k]$. Untuk motor posisi dengan $x_a = \theta$ dan $x_b = \omega$, $\Phi_{bb} = 0{,}8187$ dan $\Phi_{ab} = 0{,}09063$; kutub pengamat 0,4 memberi $L_r = (0{,}8187 - 0{,}4)/0{,}09063 = 4{,}62$. Pengamat ini hanya orde satu, tetapi pengukuran $y$ dipakai langsung tanpa penyaringan sehingga derau sensor masuk ke umpan balik.

\section{Prinsip pemisahan}
Bila kontroler memakai estimasi, $u = -K\hat{\mathbf{x}}$, maka dalam koordinat $(\mathbf{x}, \tilde{\mathbf{x}})$
\begin{equation}
	\begin{bmatrix}\mathbf{x}\\ \tilde{\mathbf{x}}\end{bmatrix}[k+1] = \begin{bmatrix}\Phi - \Gamma K & \Gamma K\\ 0 & \Phi - LC\end{bmatrix}\begin{bmatrix}\mathbf{x}\\ \tilde{\mathbf{x}}\end{bmatrix}[k].
\end{equation}
Karena matriksnya segitiga blok, kutub sistem gabungan adalah gabungan nilai eigen $\Phi - \Gamma K$ dan $\Phi - LC$. Dengan demikian $K$ dan $L$ dapat dirancang terpisah. Untuk contoh di atas, kutub sistem adalah $0{,}7252 \pm j0{,}2130$, $0{,}40$, dan $0{,}45$. Aturan praktisnya, kutub pengamat dibuat 2--4 kali lebih cepat daripada kutub kontroler: terlalu cepat membuat estimasi berderau, terlalu lambat memperlambat seluruh sistem.

\section{Kompensator ekuivalen}
Pengamat dan umpan balik keadaan bersama-sama membentuk satu sistem dinamis dari $y$ ke $u$:
\begin{equation}
	D(z) = \frac{U(z)}{-Y(z)} = K(zI - \Phi + \Gamma K + LC)^{-1}L.
\end{equation}
Untuk contoh (pengamat prediktor),
\[
	D(z) = \frac{9{,}336\,z - 7{,}135}{z^2 - 0{,}4816\,z + 0{,}2019},
\]
dengan zero di 0,764 dan kutub kompleks berjari-jari 0,449. Kompensator ini berperilaku seperti \textit{lead} orde dua. Jadi perancangan ruang keadaan menghasilkan kompensator klasik secara sistematis, dan margin kestabilannya dapat diperiksa melalui $D(z)G(z)$ seperti biasa.

\begin{lstlisting}[language=C++,caption={Pengamat \textit{current} dengan umpan balik keadaan},label={lst:b12-obs}]
// tiap T, y = posisi terukur (rad)
float thb = th_h + P12*w_h + G1*u_prev;   // prediksi
float wb  =        P22*w_h + G2*u_prev;
float ey  = y - thb;                      // koreksi
th_h = thb + Lc1*ey;
w_h  = wb  + Lc2*ey;
float u = constrain(Nbar*r - K1*th_h - K2*w_h, -UMAX, UMAX);
run_motor(u);
u_prev = u;      // gunakan u setelah saturasi
\end{lstlisting}

Gunakan $u$ setelah saturasi untuk pengamat agar estimasi tetap konsisten saat aktuator jenuh, inisialisasi $\hat\theta$ dengan pengukuran pertama, dan pastikan model yang dipakai sama dengan model perancangan.

\begin{ringkasan}
	Pengamat prediktor memiliki galat $\Phi - LC$; pengamat \textit{current} memiliki galat $\Phi - L_cC\Phi$. $L$ dirancang secara dual dengan penempatan kutub. Pengamat orde tereduksi hanya mengestimasi keadaan yang tidak terukur. Prinsip pemisahan: kutub sistem adalah gabungan kutub kontroler dan kutub pengamat.\\[3pt]
	\en{Observers reconstruct unmeasured states from model, input, and output; observer gains follow from dual pole placement; reduced-order observers estimate only unmeasured states; the separation principle allows independent design of $K$ and $L$.}
\end{ringkasan}

\paragraph{Istilah kunci.} pengamat \en{(observer)}, estimasi keadaan \en{(state estimation)}, pengamat prediktor/\textit{current} \en{(prediction/current observer)}, orde tereduksi \en{(reduced-order)}, dualitas \en{(duality)}, prinsip pemisahan \en{(separation principle)}.

\section*{Latihan}
\begin{enumerate}
	\item Hitung $L$ pengamat prediktor untuk kutub 0,5 dan 0,55, lalu bandingkan respons galatnya dengan kutub 0,40 dan 0,45.
	\item Turunkan dinamika galat pengamat \textit{current}.
	\item Rancang pengamat orde tereduksi dengan kutub 0,3 dan simulasikan dengan derau pengukuran.
	\item Hitung $D(z)$ untuk kombinasi $K$ (Bab~\ref{bab:pp}) dan $L_c$, lalu periksa margin fasanya.
	\item Mengapa pengamat dengan kutub di $z = 0$ (\textit{deadbeat}) peka terhadap derau pengukuran?
\end{enumerate}
